\documentclass{article}
\usepackage[UTF8]{ctex}
\usepackage{amsmath,amssymb,amsfonts}
\usepackage{enumitem}
\usepackage{geometry}
\geometry{a4paper, margin=1in}
% === 错题本排版优化（不参与题目标识） ===
\setlength{\parindent}{0pt}
\setlength{\parskip}{0.35em plus 0.10em minus 0.08em}
\setlength{\emergencystretch}{3em}
\setlist{itemsep=0.25em, topsep=0.40em, parsep=0pt}
\allowdisplaybreaks[2]
% === 错题本排版优化结束 ===
\begin{document}
\section*{四川大学2024年高等代数真题}
\section*{155．四川大学2024年高等代数试卷（回忆版）}

\subsection{题目1}
一。设 \(\alpha=\sqrt{p}+\sqrt{q}\) ，其中 \(p\) 是素数，\(q\) 是正整数，且 \(q \nmid q\) ．

\begin{itemize}
  \item[（1）] 证明：\(\alpha\) 是无理数。
  \item[（2）] 设 \(W=\{\beta=h(\alpha) \mid h(x) \in \mathbb{R}[x]\}\) ，证明： \(\operatorname{dim} W \leq 4\) ．
\end{itemize}

\subsection{题目2}
二．设 \(f(x)=x^{2024}+6 x-3, S\) 是 \(f(x)\) 所有复根的集合．

\begin{itemize}
  \item[（1）] 求 \(S\) 中元的个数．
  \item[（2）] 若对 \(n\) 阶矩阵 \(A\) ，有 \(f(A)=O\) ，写出所有满足 \(h(A)=O\) 的 \(h(x)\) ．
\end{itemize}
\end{document}
